\paraid{p-d1bcda9ea36a}
We now approximate the distance functions on a compact metric space by
functions in a finite-dimensional vector space.
The selected probability measure makes this vector space invariant under
isometries.
Using the spectral results from Section~\ref{sec:prelim-spectral},
we first prove approximation at one space.
We then compare the resulting functions and coordinates on nearby spaces.

\paraid{p-a4c5f187e680}
A simple construction explains why finite-dimensional coordinates can
approximate a metric.
Let
$\MetricSpace{\BaseCarrier}{\MetricSymbol}$
be nonempty and compact,
fix
$\ErrorTolerance>0$,
and choose a finite
$\ErrorTolerance$-net
$\{\NetPoint_0,\ldots,\NetPoint_{\NetSize-1}\}$.
Define
\[
 \NetCoordinateMap\colon\BaseCarrier\to\RealNumbers^{\NetSize}
\]
by
\[
 \NetCoordinateMap(\BasePoint)
 =\bigl(\MetricSymbol(\BasePoint,\NetPoint_0),\ldots,
        \MetricSymbol(\BasePoint,\NetPoint_{\NetSize-1})\bigr).
\]
For all
$\BasePoint,\ComparisonPoint\in\BaseCarrier$,
\begin{equation}\label{eq:finite-net-distance-approximation}
 \MetricSymbol(\BasePoint,\ComparisonPoint)-2\ErrorTolerance
 \leq\norm{\NetCoordinateMap(\BasePoint)-\NetCoordinateMap(\ComparisonPoint)}_\infty
 \leq\MetricSymbol(\BasePoint,\ComparisonPoint).
\end{equation}

\paraid{p-37280d9fcd1a}
Finite families of functions also underlie Shioya's approximation of metric
measure spaces
\cite[author's manuscript, Section~4.4, Definition~4.37, Theorem~4.43 and Corollary~4.44]{ShioyaMetricMeasure}.
More precisely,
for every mm-space
$(X,d_X,\mu_X)$
and every
$\ErrorTolerance>0$,
there is an integer
$N\geq1$
for which we can choose real
$1$-Lipschitz
functions
$f_1,\ldots,f_N$
on
$X$
such that the map
$\Phi=(f_1,\ldots,f_N)$
satisfies
\[
 \BoxDistance\bigl(X,(\RealNumbers^N,\norm{\cdot}_\infty,
                         \Phi_*\mu_X)\bigr)<\ErrorTolerance.
\]
Here
$\BoxDistance$
is the distance in \Cref{def:box-distance}.
Estimate \eqref{eq:finite-net-distance-approximation}
concerns uniform error over all pairs of points
of one compact space.
Our construction also needs continuity throughout a Gromov--Hausdorff neighborhood
and compatibility with isometries.
A finite net in one space does not determine such continuous choices.
Finite spectral subspaces provide the coordinate classes of \Cref{thm:local-models}.
For a convergent sequence of spaces,
\ref{item:model-continuity} guarantees coordinate pairs with the stated
convergence of norms and coordinate maps.
Condition \ref{item:model-equivariance} specifies their transformation under isometries.
Their coordinates are the coefficients of unique best approximations.
For each compact metric space
$\MetricSpace{\ComparisonCarrier}{\ComparisonMetric}$,
we define a norm
$\NormSymbol_\ComparisonCarrier$
on the coordinate space.
Under a change of orthonormal basis,
each coordinate vector $v$ is replaced by
$\OrthogonalChange v$.
The norm in the new coordinates is therefore
$\NormSymbol_\ComparisonCarrier\circ\OrthogonalChange^{-1}$,
since a vector $v$ in the new coordinates represents
$\OrthogonalChange^{-1}v$
in the old coordinates.
Thus the distances between coordinate vectors do not depend on the basis.

\paraid{p-1bee81df97ab}
Throughout this section,
$\SelectedMeasure{\BaseCarrier}{\MetricSymbol}$
denotes the assignment of \Cref{thm:invariant-measures}.
The two subscripts denote the carrier
$\BaseCarrier$
and its metric
$\MetricSymbol$.
